\capitulo{CAPÍTULO 2}{Nociones generales del derecho civil}

\section{2.1. La responsabilidad extracontractual como fuente de obligaciones (art. 1494 del C.C.)}
\begin{cita}
\textit{«ARTÍCULO 1494. Fuente de las obligaciones. Las obligaciones nacen, ya del concurso real de las voluntades de dos o más personas, como en los contratos o convenciones; ya de un hecho voluntario de la persona que se obliga, como en la aceptación de una herencia o legado y en todos los cuasicontratos; ya a consecuencia de un hecho que ha inferido injuria o daño a otra persona, como en los delitos; ya por disposición de la ley, como entre los padres y los hijos de familia.» (Código Civil)}
\end{cita}
\textbf{Importancia.} Este artículo permite ubicar la responsabilidad extracontractual como fuente de obligaciones: indica que el delito es un hecho jurídico y que los hechos jurídicos son fuente de la responsabilidad extracontractual (pueden generar obligaciones de reparar). Además, muestra:

\begin{itemize}
\item La autonomía, pero también las relaciones, entre la responsabilidad extracontractual y la penal.
\item La autonomía entre la responsabilidad extracontractual y la contractual.
\end{itemize}
\begin{comentario}
\textcolor[HTML]{8A6D3B}{\textbf{Comentario. }}La responsabilidad extracontractual da lugar a obligaciones de carácter personal y no de carácter real. Esta precisión es la base del requisito del «carácter personal del perjuicio» (ver Capítulo 4).
\end{comentario}
\section{2.2. Fundamento de la responsabilidad extracontractual (art. 2341 del C.C.)}
\begin{cita}
\textit{«ARTÍCULO 2341. Responsabilidad extracontractual. El que ha cometido un delito o culpa, que ha inferido daño a otro, es obligado a la indemnización, sin perjuicio de la pena principal que la ley imponga por la culpa o el delito cometido.» (Código Civil)}
\end{cita}
\textbf{Importancia.} Establece el fundamento general de la responsabilidad extracontractual civil: la culpa. «Fundamento» significa el porqué se responde: en materia civil, por regla general, se responde por culpa porque el legislador así lo dispuso en la norma, y por eso debe probarse la culpa. El art. 2341 es la norma que establece el régimen subjetivo de responsabilidad.

\begin{itemize}
\item Se evoluciona de la \textit{iniuria} a la culpa como fundamento de la responsabilidad.
\item La culpa es el fundamento de la responsabilidad extracontractual civil.
\end{itemize}
\section{2.3. Noción de culpa en el Código Civil}
\begin{itemize}
\item \textbf{¿El Código Civil define la culpa extracontractual?} No. El art. 63 del C.C. define la culpa en materia contractual (manejo de negocios) y regula sus grados. No existe una definición legal de la culpa extracontractual.
\item El art. 2341 la exige como fundamento general; su contenido lo han desarrollado la doctrina y la jurisprudencia como la infracción del deber de cuidado exigible.
\item \textbf{La tridivisión del art. 63 no aplica en la responsabilidad extracontractual.} Los grados de culpa están encaminados al mundo de los contratos. En el mundo de los hechos jurídicos la tridivisión no tiene razón de ser: no importa si la culpa es leve, grave o levísima; si el daño es imputable, se responde por todo el daño, porque la reparación es integral.
\end{itemize}
\begin{definicion}
\textcolor[HTML]{1F3864}{\textbf{Definición. }}\textbf{Culpa (definición doctrinal de los hermanos Mazeaud): }«un error de conducta tal, que no lo habría cometido una persona cuidadosa situada en las mismas circunstancias externas que el autor del daño».
\end{definicion}
\begin{itemize}
\item \textbf{Análisis en abstracto:} la culpa se valora comparando la conducta del autor con la de esa persona cuidadosa.
\item \textbf{Culpa y culpabilidad no son sinónimos.} La acción de repetición, por ejemplo, procede por dolo o culpa grave y exige un juicio específico de culpabilidad del agente (ver Capítulo 10).
\end{itemize}
\section{2.4. Regímenes subjetivos de responsabilidad}
\begin{tabla}
\begin{longtable}{|>{\raggedright\arraybackslash}m{\dimexpr 0.394\linewidth-2\tabcolsep-0.600pt\relax}|>{\raggedright\arraybackslash}m{\dimexpr 0.606\linewidth-2\tabcolsep-0.600pt\relax}|}
\hline
\cellcolor[HTML]{1F3864}\textcolor[HTML]{FFFFFF}{\textbf{Elementos}} & \cellcolor[HTML]{1F3864}\textcolor[HTML]{FFFFFF}{\textbf{Exoneración}} \\ \hline
1. Daño o perjuicio (para los administrativistas son conceptos distintos) & Ausencia de daño o perjuicio \\ \hline
\cellcolor[HTML]{F2F5FA}2. Nexo de causalidad: demostrar que el daño proviene de una acción u omisión del demandado & \cellcolor[HTML]{F2F5FA}Ruptura del nexo causal, probando alguna de estas causas: hecho exclusivo de un tercero, hecho exclusivo de la víctima, fuerza mayor o caso fortuito \\ \hline
3. Culpa: debe acreditarse para que haya condena & Ausencia de culpa o diligencia: probar que se actuó de manera diligente y prudente \\ \hline
\end{longtable}
\end{tabla}
El demandado se exonera destruyendo alguno de los tres elementos.

\subsection{2.4.1. Aspecto probatorio: el art. 167 del CGP}
A cada parte le corresponde probar los hechos en que funda sus afirmaciones y excepciones. Si se es demandante, se debe probar el daño, el nexo y la culpa; si el demandado se quiere exonerar, debe probar las excepciones que alega (por ejemplo, el hecho de un tercero, o que obró con diligencia) o controvertir los hechos.

\subsection{2.4.2. La carga dinámica de la prueba}
\begin{definicion}
\textcolor[HTML]{1F3864}{\textbf{Definición. }}\textbf{Carga dinámica de la prueba: }el juez, de oficio o a petición de parte, puede distribuir la carga de probar determinados hechos y asignarla a quien se encuentre en mejores condiciones para acreditarlos, por tener mayor cercanía, disponibilidad o facilidad de acceso a la prueba.
\end{definicion}
\begin{itemize}
\item Puede hacerse hasta antes de dictar sentencia.
\item No es una inversión automática de toda la carga de la prueba, sino una redistribución según las circunstancias concretas del caso.
\item Es frecuente en responsabilidad médica, donde el profesional o la institución de salud suele estar en una posición probatoria más favorable respecto de aspectos técnicos o científicos.
\end{itemize}
\section{2.5. La revolución industrial y los regímenes objetivos}
\subsection{2.5.1. Origen histórico: influencia de la revolución industrial}
\begin{itemize}
\item Surgimiento de actividades peligrosas en las fábricas y en la sociedad (máquinas, vehículos).
\item Necesidad de proteger a los trabajadores de eventuales riesgos en el trabajo.
\item Necesidad de proteger a los consumidores de los riesgos creados por actividades peligrosas.
\end{itemize}
\subsection{2.5.2. Consecuencias: el favorecimiento de la víctima}
\begin{itemize}
\item Régimen objetivo de responsabilidad para actividades peligrosas, originado a partir de una interpretación jurisprudencial «amplia» del art. 1384 del Código Civil francés (responsabilidad por el hecho de las cosas).
\item Leyes relativas a accidentes de trabajo e indemnizaciones tarifadas.
\item Socialización de riesgos a través de contratos de seguros obligatorios y no obligatorios.
\end{itemize}
Con la revolución industrial aumentó el uso de máquinas, vehículos y otros instrumentos capaces de generar riesgos significativos. Esto llevó a replantear el régimen tradicional de responsabilidad extracontractual, fundado principalmente en la culpa. Surgieron regímenes de responsabilidad objetiva, en los que la obligación de reparar no depende necesariamente de demostrar culpa, sino de la creación o materialización de un riesgo. Este desarrollo tuvo una importante influencia de la jurisprudencia francesa. Paralelamente, el crecimiento de estas actividades riesgosas impulsó la utilización de contratos de seguro como mecanismo para distribuir y cubrir económicamente los riesgos.

\subsection{2.5.3. Responsabilidad extracontractual por actividades peligrosas}
\begin{itemize}
\item \textbf{Fundamento: la creación del riesgo.} No está principalmente en la culpa, sino en el riesgo creado por el desarrollo de una actividad peligrosa o por el uso de una cosa peligrosa.
\item \textbf{Responsabilidad del guardián.} Responde quien tiene la guarda material o jurídica de la actividad o del instrumento, es decir, quien ejerce poder de dirección, control y manejo sobre este y obtiene provecho de su utilización.
\end{itemize}
\begin{ejemplo}
\textcolor[HTML]{4F6228}{\textbf{Ejemplo. }}Si Óscar es propietario y guardián de una motocicleta y conserva sobre ella poder de dirección y control, puede llegar a responder por los daños causados con su utilización, incluso cuando sea otra persona quien la conduzca, dependiendo de las circunstancias del caso.
\end{ejemplo}
\begin{itemize}
\item \textbf{Clases de cosas peligrosas:}
\begin{itemize}
\item Por su estructura: cuando su propia naturaleza implica un riesgo especial.
\item Por su comportamiento: cuando adquieren peligrosidad por la forma en que son utilizadas, manipuladas o puestas en funcionamiento.
\end{itemize}
\end{itemize}
\subsection{2.5.4. Regímenes objetivos de responsabilidad: elementos y exoneración}
\begin{tabla}
\begin{longtable}{|>{\raggedright\arraybackslash}m{\dimexpr 0.482\linewidth-2\tabcolsep-0.600pt\relax}|>{\raggedright\arraybackslash}m{\dimexpr 0.518\linewidth-2\tabcolsep-0.600pt\relax}|}
\hline
\cellcolor[HTML]{1F3864}\textcolor[HTML]{FFFFFF}{\textbf{Elementos}} & \cellcolor[HTML]{1F3864}\textcolor[HTML]{FFFFFF}{\textbf{Exoneración}} \\ \hline
1. Daño o perjuicio & Ausencia de daño o perjuicio \\ \hline
\cellcolor[HTML]{F2F5FA}2. Nexo causal entre la actividad peligrosa y el daño & \cellcolor[HTML]{F2F5FA}Ruptura del nexo de causalidad mediante una causa extraña \\ \hline
\end{longtable}
\end{tabla}
La diligencia no exonera, porque está en el ámbito de la culpa y la culpa no es elemento del régimen: el demandado no puede decir que iba manejando bien, que paró o que fue diligente.

\begin{comentario}
\textcolor[HTML]{8A6D3B}{\textbf{Comentario. }}Conclusión. En responsabilidad extracontractual la regla general es el régimen subjetivo, fundado en la culpa. Sin embargo, existen eventos de régimen objetivo, especialmente cuando el daño proviene del desarrollo de una actividad peligrosa o del uso de una cosa o instrumento peligroso; allí el fundamento se desplaza de la culpa al riesgo creado.
\par
\textcolor[HTML]{8A6D3B}{\textbf{Comentario. }}Síntesis. En materia civil, el fundamento general de la responsabilidad es la culpa (art. 2341 del C.C.). Excepcionalmente, el fundamento es el riesgo: cuando se trata de actividades peligrosas o instrumentos peligrosos, lo que permite aplicar un régimen objetivo. En el ámbito contractual, la responsabilidad objetiva se aplica cuando hay obligaciones de custodia u obligaciones de resultado. Lo general es el régimen subjetivo y, excepcionalmente, en estas pocas situaciones, el régimen objetivo. La lógica en la responsabilidad estatal es bastante parecida (ver Capítulo 6).
\end{comentario}
\section{2.6. Responsabilidad directa e indirecta}
\begin{definicion}
\textcolor[HTML]{1F3864}{\textbf{Definición. }}\textbf{Responsabilidad directa (por el hecho propio): }se presenta cuando en una misma persona coinciden el autor material del daño y el responsable de repararlo; quien causa el daño es también quien debe responder por sus consecuencias.
\par
\textcolor[HTML]{1F3864}{\textbf{Definición. }}\textbf{Responsabilidad indirecta (por el hecho ajeno): }se presenta cuando el autor material del daño y la persona llamada a responder son sujetos diferentes.
\end{definicion}
\begin{tabla}
\begin{longtable}{|>{\raggedright\arraybackslash}m{\dimexpr 0.134\linewidth-2\tabcolsep-0.533pt\relax}|>{\raggedright\arraybackslash}m{\dimexpr 0.274\linewidth-2\tabcolsep-0.533pt\relax}|>{\raggedright\arraybackslash}m{\dimexpr 0.592\linewidth-2\tabcolsep-0.533pt\relax}|}
\hline
\cellcolor[HTML]{1F3864}\textcolor[HTML]{FFFFFF}{\textbf{Aspecto}} & \cellcolor[HTML]{1F3864}\textcolor[HTML]{FFFFFF}{\textbf{Responsabilidad directa}} & \cellcolor[HTML]{1F3864}\textcolor[HTML]{FFFFFF}{\textbf{Responsabilidad indirecta}} \\ \hline
Fundamento legal & Art. 2341 del C.C. & Arts. 2347 y 2349 del C.C. \\ \hline
\cellcolor[HTML]{F2F5FA}Fundamento de la imputación & \cellcolor[HTML]{F2F5FA}Culpa o creación de riesgo por el causante del daño & \cellcolor[HTML]{F2F5FA}Culpa o creación de riesgo por el causante del daño, más culpa \textit{in eligendo} (elección inadecuada) o \textit{in vigilando} (incumplimiento del deber de vigilancia o control) del responsable \\ \hline
Prescripción & Ordinaria del art. 2536 del C.C. (10 años), salvo regla especial & Especial del art. 2358 del C.C. (3 años contados desde la perpetración del acto) \\ \hline
\cellcolor[HTML]{F2F5FA}Sujetos & \cellcolor[HTML]{F2F5FA}Causante del daño = responsable del daño & \cellcolor[HTML]{F2F5FA}Causante del daño ≠ responsable del daño \\ \hline
\end{longtable}
\end{tabla}
\begin{itemize}
\item El art. 2347 del C.C. establece que una persona puede responder por los hechos de quienes se encuentren bajo su cuidado o dependencia. Ejemplos tradicionales: la responsabilidad de los padres por determinados hechos de sus hijos y la del empleador por los hechos de sus dependientes, en los casos previstos por la ley.
\item Frente a la víctima puede operar una presunción de culpa indirecta o mediata, de modo que no siempre debe demostrar directamente que el responsable escogió o vigiló mal.
\item El término de tres años del art. 2358 debe distinguirse del régimen general aplicable a la responsabilidad directa.
\end{itemize}
\begin{ejemplo}
\textcolor[HTML]{4F6228}{\textbf{Ejemplo. }}Una empresa contrata a un mensajero que, durante el desarrollo de sus funciones, causa un accidente. Dependiendo de las circunstancias, la víctima puede reclamar no solo contra el autor material del daño, sino también contra quien jurídicamente deba responder por el hecho ajeno.
\par
\textcolor[HTML]{4F6228}{\textbf{Ejemplo. }}Caso: accidente de tránsito de un bus de una empresa de servicio público. La víctima no tiene que probar la culpa, porque la conducción de vehículos es una actividad peligrosa (art. 2356 del C.C.). Frente a ella responden:
\end{ejemplo}
\begin{itemize}
\item \textbf{Conductor:} responsabilidad directa por actividad peligrosa (régimen objetivo); prescripción de 10 años.
\item \textbf{Propietario del bus:} responsabilidad directa como guardián de la actividad o del instrumento, según quién conserve la dirección y el control; prescripción de 10 años.
\item \textbf{Empresa transportadora:} responsabilidad directa, como guardiana de la actividad peligrosa; prescripción de 10 años (ver la precisión siguiente).
\end{itemize}
\begin{precision}
\textcolor[HTML]{9C2B2B}{\textbf{Precisión. }}La empresa transportadora \textbf{no} responde por culpa \textit{in vigilando} o \textit{in eligendo} ni tiene el término de 3 años. Según la jurisprudencia de la Sala de Casación Civil de la Corte Suprema de Justicia: (i) la responsabilidad de las personas jurídicas es \textbf{directa}: el hecho de sus agentes, empleados o vinculados se considera hecho propio de la persona jurídica (tesis acogida desde el fallo de 1962 y reiterada desde entonces); (ii) la empresa a la que está afiliado el vehículo es \textbf{guardiana de la actividad peligrosa} de conducción y responde con fundamento en el \textbf{artículo 2356 del C.C.} (actividades peligrosas), de forma \textbf{solidaria} con el propietario y el conductor (art. 2344 del C.C.); la Corte ha dicho que el solo hecho de la afiliación implica, en principio, que la empresa soporte responsabilidad, y esa presunción solo se desvirtúa probando, por ejemplo, el hurto del vehículo o la transferencia válida de su tenencia; (iii) por ser responsabilidad directa, el término de prescripción es el \textbf{ordinario de 10 años} (arts. 2512 y 2536 del C.C., este último modificado por la Ley 791 de 2002), no el de 3 años del inciso 2 del art. 2358, que solo aplica a los terceros que responden por el hecho ajeno (por ejemplo, ciertos casos de padres o empleadores personas naturales). En síntesis: conductor, propietario y empresa responden los tres de forma directa, bajo el régimen de actividades peligrosas, solidariamente y con prescripción de 10 años; en ese régimen la víctima no tiene que probar la culpa, y la diligencia no exonera: solo la causa extraña.
\end{precision}
\section{2.7. Relaciones entre la responsabilidad extracontractual y la penal}
\textbf{El delito como fuente de daños.} Un mismo hecho puede dar lugar a dos acciones: la acción penal y la acción civil. La absolución penal no siempre impide una condena civil.

\begin{cita}
\textit{«ARTÍCULO 2358. Prescripción de la acción de reparación. Las acciones para la reparación del daño proveniente de delito o culpa que puedan ejercitarse contra los que sean punibles por el delito o la culpa, se prescriben dentro de los términos señalados en el Código Penal para la prescripción de la pena principal. Las acciones para la reparación del daño que puedan ejercitarse contra terceros responsables, conforme a las disposiciones de este capítulo, prescriben en tres años contados desde la perpetración del acto.» (Código Civil)}
\end{cita}
\textbf{Caminos de la víctima:}

\begin{itemize}
\item Si la acción penal prescribe, se pierde el incidente de reparación integral de perjuicios.
\item Se puede adelantar al mismo tiempo el proceso penal y el proceso civil; la pregunta es hasta qué punto la decisión penal ata la decisión del juez civil.
\end{itemize}
\textbf{Elemento en común: la causalidad (Ley 600 de 2000).} Existían dos artículos en la Ley 600 de 2000 que regulaban los efectos de la cosa juzgada penal en el proceso de reparación de perjuicios. La Ley 600 fue sustituida por la Ley 906 de 2004 (sistema penal acusatorio), que no reprodujo estas reglas; la Ley 600 solo sigue aplicándose a los procesos por hechos ocurridos antes de la entrada en vigencia de la Ley 906:

\begin{cita}
\textit{«Art. 57. Efectos de la cosa juzgada penal absolutoria. La acción civil no podrá iniciarse ni proseguirse cuando se haya declarado, por providencia en firme, que la conducta causante del perjuicio no se realizó o que el sindicado no lo cometió o que obró en estricto cumplimiento de un deber legal o en legítima defensa.»}
\par
\textit{«Art. 59. Efectos de la declaratoria de responsabilidad. Cuando no se hubiere constituido parte civil y se condene al procesado, la responsabilidad no podrá ser discutida en el proceso civil, debiendo limitarse éste a la clase y monto de los perjuicios.»}
\end{cita}
\textbf{Las cuatro razones del art. 57.} Si la absolución se daba por una de las cuatro razones, el proceso civil debía terminar automáticamente; si no, continuaba:

\begin{itemize}
\item \textbf{Razón 1 – La conducta no se realizó:} no hubo daño (falta el hecho dañoso).
\item \textbf{Razón 2 – El sindicado no la cometió:} no hubo nexo causal.
\item \textbf{Razón 3 – Obró en estricto cumplimiento de un deber legal:} solo importa, y solo exonera, si se está aplicando un régimen subjetivo de responsabilidad extracontractual.
\item \textbf{Razón 4 – Obró en legítima defensa:} rompe el nexo causal como hecho exclusivo de la víctima (ver la precisión siguiente).
\end{itemize}
\begin{precision}
\textcolor[HTML]{9C2B2B}{\textbf{Precisión. }}La legítima defensa no es un «hecho de un tercero». Cuando quien sufre el daño es el propio agresor, la legítima defensa rompe el nexo causal como \textbf{hecho exclusivo (y determinante) de la víctima}: fue la agresión injusta de la víctima la que causó el daño. Así la trata el Consejo de Estado cuando agentes del Estado alegan haber actuado en legítima defensa: la ubica entre las causales eximentes (hecho de un tercero, hecho de la víctima, fuerza mayor) como manifestación del hecho exclusivo de la víctima, y exige que el Estado la pruebe y que el uso de la fuerza haya sido proporcional y razonado (matar es la \textit{ultima ratio}); si no se prueba la agresión o la reacción fue desproporcionada, no hay exoneración. En consecuencia, solo exonera si se prueba una defensa necesaria y proporcional frente a una agresión actual e injusta. Si el daño lo sufre un tercero ajeno a la agresión (por ejemplo, un transeúnte herido por la bala del defensor), la legítima defensa no le es oponible como hecho de ese tercero, y en el caso del Estado puede haber responsabilidad objetiva (ver Capítulo 6, enfrentamientos armados).
\end{precision}
En cualquier otro caso de absolución, el juez civil puede condenar: puede subsistir responsabilidad civil, porque los presupuestos y los estándares probatorios de ambos procesos son diferentes. La lógica de estos cuatro casos se sigue aplicando hoy, a pesar de que la norma no está en la Ley 906.

\textbf{Art. 59.} Con la sentencia penal condenatoria quedaba acreditada la responsabilidad, y en el proceso civil solo se discutía la tasación (clase y monto) del perjuicio, no el nexo causal.

\begin{comentario}
\textcolor[HTML]{8A6D3B}{\textbf{Comentario. }}Conclusión: aunque las dos normas no están vigentes, su lógica se sigue usando y debe explicársele al juez.
\end{comentario}
\section{2.8. Funciones de la responsabilidad extracontractual}
\textbf{1. Reparatoria o compensatoria:} busca reparar integralmente el daño causado.

\textbf{2. Preventiva:} parte de la doctrina se la atribuye, pero es discutible: el proceso no busca prevenir el daño, porque cuando se demanda el daño ya fue causado. Las acciones preventivas por excelencia son, por ejemplo, la acción popular y la tutela.

\textbf{3. Demarcatoria:} busca delimitar qué conductas son jurídicamente admisibles y cuáles no dentro de la sociedad, estableciendo estándares de comportamiento y señalando cuándo una actuación que causa daño genera el deber de responder. Se predica principalmente de los regímenes subjetivos, porque en ellos se valora la conducta del agente a partir de la culpa o el dolo; en los regímenes objetivos la responsabilidad no depende de demostrar que la conducta fue reprochable: no hay lugar a demarcar.

\textbf{4. Distributiva:} busca distribuir socialmente las consecuencias económicas del daño, especialmente a través de mecanismos de aseguramiento obligatorio o voluntario. El costo del daño no recae exclusivamente sobre quien lo causa, sino que se reparte entre un conjunto de personas que contribuyen al sistema mediante primas o aportes.

\begin{ejemplo}
\textcolor[HTML]{4F6228}{\textbf{Ejemplo. }}El SOAT: los propietarios de vehículos pagan una prima y, cuando ocurre un accidente de tránsito, el seguro cubre determinadas prestaciones con cargo al fondo común formado por esas contribuciones. Así, el riesgo y sus consecuencias económicas se redistribuyen entre todos los asegurados.
\end{ejemplo}
\textbf{5. Sancionatoria:} en principio, la responsabilidad extracontractual no busca sancionar al responsable, sino reparar integralmente el daño. En Estados Unidos existen los daños punitivos, que sí cumplen una función sancionatoria: el juez puede imponer una suma superior al valor del perjuicio para castigar conductas especialmente graves y disuadir su repetición; su cuantía puede atender al grado de reproche o culpabilidad de la conducta. En Colombia, como regla general, no existe la figura autónoma del daño punitivo en la responsabilidad civil: la indemnización debe corresponder al perjuicio efectivamente causado.

\begin{comentario}
\textcolor[HTML]{8A6D3B}{\textbf{Comentario. }}En principio la indemnización llega hasta el monto del perjuicio: «todo el perjuicio, nada más y nada menos que el perjuicio». Con esa base puede cuestionarse si la posibilidad de triplicar los perjuicios morales en graves violaciones de derechos humanos no introduce una función sancionatoria (ver Capítulo 4, perjuicios morales).
\end{comentario}
\textbf{6. Expresiva:} busca transmitir un mensaje de reconocimiento, reivindicación y restablecimiento de los derechos vulnerados, mediante medidas que no son exclusivamente económicas: medidas de satisfacción como disculpas públicas, actos de reconocimiento de responsabilidad, rectificaciones o publicación de decisiones. En Colombia, el Consejo de Estado ha ordenado este tipo de medidas en casos de graves vulneraciones de derechos, incluida la privación injusta de la libertad, para reconocer públicamente la afectación sufrida por la víctima y contribuir a la reparación integral.

\section{2.9. ¿Quién define el fundamento de la responsabilidad y sus consecuencias?}
\begin{itemize}
\item La visión de la responsabilidad puede construirse desde la víctima o desde el responsable.
\item ¿Quién lo define? El legislador.
\item \textbf{Hoy, en la práctica,} es el juez quien define si el régimen es subjetivo u objetivo, porque ni la ley ni la Constitución regulan el tema; por eso, por los mismos hechos, se falla con regímenes diferentes. Cuando la ley o la Constitución no fijan expresamente un único régimen, el juez determina el título de imputación aplicable según los hechos y la jurisprudencia.
\end{itemize}
\subsection{2.9.1. ¿A quién favorece cada régimen?}
\begin{itemize}
\item \textbf{Régimen subjetivo: favorece más al demandado} en términos probatorios, porque impone al demandante una carga adicional: no basta demostrar el daño y el nexo causal; también debe acreditar la culpa o el dolo, salvo que opere alguna presunción. Eso hace más difícil estructurar la responsabilidad.
\item \textbf{Régimen objetivo: favorece más a la víctima (demandante),} porque no tiene que probar la culpa o el dolo; le basta acreditar los elementos de ese régimen, principalmente el daño, el nexo causal y la imputación. El demandado no se exonera demostrando que actuó con diligencia: normalmente debe acreditar una causa extraña que rompa el nexo causal (fuerza mayor, hecho exclusivo de la víctima o hecho exclusivo de un tercero).
\end{itemize}
\subsection{2.9.2. La responsabilidad médica}
La regla general en Colombia es que la responsabilidad médica es subjetiva: no basta con que el paciente haya sufrido un daño; debe demostrarse una falla en la conducta médica (negligencia, imprudencia, impericia o incumplimiento de la \textit{lex artis}), además del daño y del nexo causal. El médico no responde simplemente porque el resultado haya sido desfavorable. Sin embargo, la carga de la prueba puede distribuirse dinámicamente cuando la institución o el profesional están en mejores condiciones de acreditar hechos técnicos. Hay situaciones específicas en que ciertas obligaciones pueden aproximarse a obligaciones de resultado, pero la atención médica ordinaria es, por regla general, una obligación de medio. (Sobre las excepciones en la responsabilidad médica del Estado —ginecobstetricia, cirugías estéticas, infecciones nosocomiales—, ver Capítulo 6.)
